\section*{Anexo B. Cuadros resumen de varias vidas}
\label{anexo:resumen-b}

\textbf{Población y asignación inicial.} Los cuadros resumen la población
de A utilizada en B y las parejas hipotéticas asignadas. Los menores de
30 años permanecen en el seguro colectivo, pero no generan pagos en B.

\begin{table}[htbp]
\centering
\small
\begin{tabular}{@{}lr@{}}
\toprule
Concepto & Cantidad \\
\midrule
Activos iniciales A & 104\,004 \\
Sin pareja: menores de 30 & 6\,755 \\
Con pareja teorica & 97\,249 \\
Exclusiones propias de B & 0 \\
Parejas admisibles B & 97\,249 \\
\bottomrule
\end{tabular}

\caption{Población inicial y admisibilidad de parejas en B. Unidad: personas
o parejas, según el concepto. Elaboración propia.}
\label{tab:b-poblacion}
\end{table}

\begin{table}[htbp]
\centering
\small
\begin{tabular}{@{}llrrr@{}}
\toprule
Edad inicial activo & Sexo activo & $\delta$ & Edad pareja & Parejas admisibles \\
\midrule
30 a 34 & Mujer & 0 & 30--34 & 5\,067 \\
30 a 34 & Hombre & 0 & 30--34 & 3\,730 \\
35 a 54 & Mujer & 3 & 38--57 & 24\,300 \\
35 a 54 & Hombre & 3 & 32--51 & 21\,488 \\
55 o mas & Mujer & 5 & 60--101 & 19\,501 \\
55 o mas & Hombre & 5 & 50--91 & 23\,163 \\
\bottomrule
\end{tabular}

\caption{Parejas admisibles por banda de edad y sexo del activo.
La edad de la pareja muestra el intervalo observado al inicio. Elaboración propia.}
\label{tab:b-bandas}
\end{table}

Cada pareja tiene sexo opuesto al activo. Para un hombre se aplica
$y=x-\delta$ y para una mujer $y=x+\delta$. La diferencia de edad se fija
al inicio y se conserva durante los cinco años.

\textbf{Probabilidades de pago.} Para resumir las cohortes se calcula, en
cada cierre, la media ponderada por las cantidades iniciales admisibles:
$\overline{P}_h(t)=\sum n_{x,s}P_h(x,s;t)/\NParejasB$.
Estas probabilidades usan mortalidad biológica, sin renuncia.

\begin{table}[htbp]
\centering
\small
\begin{tabular}{@{}rrrr@{}}
\toprule
Año & Ambos vivos & Último sobreviviente & Reversión \\
\midrule
1 & 97,94\% & 99,96\% & 0,96\% \\
2 & 95,90\% & 99,83\% & 1,88\% \\
3 & 93,86\% & 99,62\% & 2,77\% \\
4 & 91,82\% & 99,33\% & 3,62\% \\
5 & 89,78\% & 98,97\% & 4,45\% \\
\bottomrule
\end{tabular}

\caption{Probabilidad media de pago al cierre de cada año, ponderada
por las parejas iniciales. Elaboración propia.}
\label{tab:b-probabilidades}
\end{table}

Último sobreviviente significa al menos una persona viva, incluidas ambas.
Reversión significa cónyuge vivo con el activo ya fallecido.

\clearpage
\subsection*{Valores presentes de las tres rentas}

Cada renta paga una unidad al cierre de cada año en que se cumple su
estatus, durante cinco años, con descuento efectivo del $4\%$.
Los VP se expresan en unidades de renta descontadas y corresponden a
tres coberturas alternativas.

\begin{table}[htbp]
\centering
\small
\begin{tabular}{@{}lr@{}}
\toprule
Estatus de pago & VP agregado \\
\midrule
Ambos vivos & 407\,048,93 \\
Al menos uno vivo (ultimo sobreviviente) & 431\,032,88 \\
Conyuge vivo con activo fallecido & 11\,555,86 \\
\bottomrule
\end{tabular}

\caption{VP agregados para \NParejasB{} parejas admisibles.
Elaboración propia.}
\label{tab:vp-varias-vidas}
\end{table}

\begin{table}[htbp]
\centering
\small
\begin{tabular}{@{}lrrrr@{}}
\toprule
Sexo del activo & Parejas & VP ambos & VP último & VP reversión \\
\midrule
Mujer & 48\,868 & 203\,197,12 & 216\,296,65 & 3\,787,98 \\
Hombre & 48\,381 & 203\,851,81 & 214\,736,23 & 7\,767,88 \\
\bottomrule
\end{tabular}

\caption{Descomposición de los VP por sexo del activo.
Las dos filas suman el VP agregado de cada cobertura antes del redondeo.
Elaboración propia.}
\label{tab:b-vp-sexo}
\end{table}

\textbf{Ambos vivos.} Su VP es \VPambosB{}. Los pagos requieren la
supervivencia simultánea de las dos personas y cesan tras el primer
fallecimiento. Por ello presenta un valor inferior al de último sobreviviente.

\textbf{Último sobreviviente.} Su VP es \VPultimoB{}, el mayor de las
tres alternativas. La cobertura sigue pagando mientras viva cualquiera
de las dos personas, incluso si el cónyuge falleció y solo vive el activo.

\textbf{Reversión al cónyuge.} Su VP es \VPreversionB{}. El pago se
activa cuando el activo ya falleció y su pareja sigue viva, y puede
repetirse en cierres posteriores. Es distinto de una indemnización
única por fallecimiento.

Los tres VP no se suman entre sí ni se incorporan al seguro colectivo.
Para una renta anual de importe $a$, cada VP se multiplica por $a$.
La descomposición por sexo también refleja la composición de edades
y la regla de asignación de parejas de cada grupo.
